% !TEX program = pdflatex
% ==========================================================
% CLASE BEAMER MODERNA Y ESTRUCTURADA
% Fundamentos de Matemática Básica - Agronomía
% Tema: Potencias
% Universidad Santo Tomás - Talca
% ==========================================================

\newif\ifhandout
%\handoutfalse
\handouttrue

\ifhandout
\documentclass[aspectratio=169,10pt,handout]{beamer}
\else
\documentclass[aspectratio=169,10pt]{beamer}
\fi

% ----------------------------------------------------------
% PAQUETES
% ----------------------------------------------------------
\usepackage[spanish,es-noshorthands]{babel}
\usepackage[T1]{fontenc}
\usepackage[utf8]{inputenc}
\usepackage{lmodern}
\usepackage{amsmath,amsfonts,amssymb}
\usepackage{ragged2e}
\usepackage{enumitem}
\usepackage{multicol}
\usepackage{tikz}
\usetikzlibrary{positioning}
\usepackage{fontawesome5}

% ----------------------------------------------------------
% COLORES
% ----------------------------------------------------------
\definecolor{USTBlue}{RGB}{0,70,127}
\definecolor{USTBlueDark}{RGB}{0,45,85}
\definecolor{USTTeal}{RGB}{0,153,117}
\definecolor{USTGray}{RGB}{120,120,120}
\definecolor{USTSoft}{RGB}{248,249,250}
\definecolor{USTText}{RGB}{35,40,45}
\definecolor{WarmGray}{RGB}{120,120,120}

% ----------------------------------------------------------
% TEMA
% ----------------------------------------------------------
\usetheme{Madrid}
\usecolortheme[named=USTBlue]{structure}
\setbeamertemplate{navigation symbols}{}

% ----------------------------------------------------------
% PIE DE PÁGINA
% ----------------------------------------------------------
\setbeamercolor{author in head/foot}{bg=USTSoft,fg=USTGray}
\setbeamercolor{date in head/foot}{bg=USTSoft,fg=USTGray}
\setbeamercolor{title in head/foot}{bg=USTSoft,fg=USTGray}

\setbeamertemplate{footline}{%
	\leavevmode%
	\hbox{%
		\begin{beamercolorbox}[wd=.85\paperwidth,ht=2.5ex,dp=1ex,leftskip=0.4cm]{author in head/foot}%
			\scriptsize
			\textbf{Fundamentos de Matemática Básica}
			\quad|\quad
			Agronomía
			\quad|\quad
			Universidad Santo Tomás Talca
		\end{beamercolorbox}%
		\begin{beamercolorbox}[wd=.15\paperwidth,ht=2.5ex,dp=1ex,rightskip=0.3cm plus1fil]{date in head/foot}%
			\scriptsize \insertframenumber/\inserttotalframenumber
		\end{beamercolorbox}%
	}%
}

% ----------------------------------------------------------
% BLOQUES
% ----------------------------------------------------------
\setbeamercolor{block title}{bg=USTBlue!15,fg=USTBlueDark}
\setbeamercolor{block body}{bg=USTSoft,fg=black}

\setbeamercolor{block title example}{bg=USTTeal!20,fg=USTBlueDark}
\setbeamercolor{block body example}{bg=USTSoft,fg=black}

\setbeamercolor{block title alerted}{bg=red!15,fg=black}
\setbeamercolor{block body alerted}{bg=red!3,fg=black}

% ----------------------------------------------------------
% COMANDOS
% ----------------------------------------------------------
\newcommand{\destacar}[1]{\textcolor{USTBlueDark}{\textbf{#1}}}
\newcommand{\alerta}[1]{\textcolor{red!70!black}{\textbf{#1}}}
\newcommand{\pp}{\pause}

\newenvironment{metaBox}
{\begin{block}{Meta de aprendizaje}}
	{\end{block}}

\newenvironment{ideaBox}
{\begin{exampleblock}{Idea clave}}
	{\end{exampleblock}}

% ----------------------------------------------------------
% PORTADA
% ----------------------------------------------------------
\newcommand{\PortadaProfesionalUST}{%
	\begin{frame}[plain]
		\begin{tikzpicture}[remember picture,overlay]
			\fill[white] (current page.south west) rectangle (current page.north east);
			\fill[USTBlueDark] (current page.north west) rectangle ([yshift=-1.35cm]current page.north east);
			\fill[USTTeal] ([yshift=-1.35cm]current page.north west) rectangle ([yshift=-1.52cm]current page.north east);
			\fill[USTSoft] (current page.south west) rectangle ([xshift=2.15cm]current page.north west);
			
			\fill[USTBlue!10] ([xshift=1.05cm,yshift=-1.8cm]current page.north west) circle (0.72cm);
			\fill[USTTeal!12] ([xshift=1.65cm,yshift=-4.3cm]current page.north west) circle (0.42cm);
			\fill[USTBlue!6] ([xshift=-1.15cm,yshift=1.0cm]current page.south east) circle (1.15cm);
			
			\fill[black!6,rounded corners=7pt]
			([xshift=-4.95cm,yshift=-3.25cm]current page.center) rectangle
			([xshift=4.95cm,yshift=2.45cm]current page.center);
			
			\fill[white,rounded corners=7pt]
			([xshift=-5.1cm,yshift=-3.1cm]current page.center) rectangle
			([xshift=4.8cm,yshift=2.6cm]current page.center);
			
			\node[anchor=north east,text=white]
			at ([xshift=-0.65cm,yshift=-0.48cm]current page.north east)
			{\small\bfseries Universidad Santo Tomás\quad Ciencias Básicas Talca};
			
			\fill[USTTeal] ([xshift=-5.1cm,yshift=-3.1cm]current page.center) rectangle
			([xshift=-4.82cm,yshift=2.6cm]current page.center);
			
			\node[align=center] at ([xshift=0.15cm,yshift=-0.1cm]current page.center){%
				\begin{minipage}{0.66\paperwidth}
					\centering
					{\Large\bfseries\color{USTText}\insertauthor\par}
					\vspace{0.95cm}
					{\fontsize{22}{26}\selectfont\bfseries\color{USTBlueDark}\inserttitle\par}
					\vspace{0.28cm}
					{\large\color{USTTeal}\insertsubtitle\par}
					\vspace{0.95cm}
					{\small\color{USTGray}\insertdate\par}
				\end{minipage}
			};
			
			\draw[USTGray!35,line width=0.35pt]
			([xshift=0.95cm,yshift=0.72cm]current page.south west) --
			([xshift=-0.95cm,yshift=0.72cm]current page.south east);
			
			\node[anchor=south,text=USTGray]
			at ([yshift=0.2cm]current page.south)
			{\scriptsize
				\textbf{Fundamentos de Matemática Básica}
				\;|\;
				Agronomía
				\;|\;
				Universidad Santo Tomás Talca};
		\end{tikzpicture}
	\end{frame}
}

% ----------------------------------------------------------
% DATOS
% ----------------------------------------------------------
\title{Fundamentos de Matemática Básica}
\subtitle{\Large Clase estructurada: Potencias}
\author{Prof. Luis González Alcaino}
\institute{Universidad Santo Tomás}
\date{Año académico 2026}

% ==========================================================
% DOCUMENTO
% ==========================================================
\begin{document}
	
	\PortadaProfesionalUST
	
	% ----------------------------------------------------------
	\begin{frame}{Plan de la clase}
		\begin{block}{Estructura de la sesión (80 minutos)}
			\begin{enumerate}
				\item \destacar{Inicio:} activación de conocimientos previos.
				\item \destacar{Desarrollo:} concepto y propiedades de las potencias.
				\item \destacar{Práctica guiada:} ejemplos resueltos paso a paso.
				\item \destacar{Práctica autónoma:} ejercicios individuales.
				\item \destacar{Cierre:} síntesis y reflexión final.
			\end{enumerate}
		\end{block}
		
		\begin{exampleblock}{Meta de la sesión}
			Aplicar el concepto de potencia y sus propiedades en ejercicios matemáticos y situaciones contextualizadas en Agronomía.
		\end{exampleblock}
	\end{frame}
	
	% ----------------------------------------------------------
	\begin{frame}{Inicio: activación de conocimientos previos}
		\begin{block}{Preguntas iniciales}
			\begin{itemize}
				\item ¿Cómo escribirías \(2\cdot2\cdot2\cdot2\) de forma abreviada?
				\item ¿Qué significa el número pequeño que aparece arriba en una potencia?
				\item ¿Dónde crees que se usan potencias en problemas reales?
			\end{itemize}
		\end{block}
		
		\pp
		
		\begin{ideaBox}
			Las potencias permiten escribir de forma breve una multiplicación repetida.
		\end{ideaBox}
	\end{frame}
	
	% ----------------------------------------------------------
	\begin{frame}{Concepto de potencia}
		\begin{block}{Definición}
			Una \destacar{potencia} representa una multiplicación repetida de un mismo número:
			\[
			a^n=\underbrace{a\cdot a\cdot a\cdots a}_{n\text{ veces}}
			\]
			donde:
			\begin{itemize}
				\item \(a\) es la \destacar{base},
				\item \(n\) es el \destacar{exponente}.
			\end{itemize}
		\end{block}
		
		\pp
		
		\begin{exampleblock}{Ejemplos}
			\[
			2^4=2\cdot2\cdot2\cdot2=16
			\]
			\[
			3^3=27
			\qquad
			10^2=100
			\]
		\end{exampleblock}
	\end{frame}
	
	% ----------------------------------------------------------
	\begin{frame}{Interpretación de base y exponente}
		\begin{columns}[T,totalwidth=\textwidth]
			\begin{column}{0.48\textwidth}
				\begin{block}{Base}
					La base es el número que se repite como factor.
					
					En \(5^3\), la base es \(5\).
				\end{block}
			\end{column}
			\begin{column}{0.48\textwidth}
				\begin{exampleblock}{Exponente}
					El exponente indica cuántas veces se repite la base.
					
					En \(5^3\), el exponente es \(3\).
				\end{exampleblock}
			\end{column}
		\end{columns}
		
		\vspace{0.3cm}
		\[
		5^3 = 5\cdot5\cdot5=125
		\]
	\end{frame}
	
	% ----------------------------------------------------------
	\begin{frame}{Propiedad 1: producto de potencias de igual base}
		\begin{block}{Regla}
			\[
			a^m\cdot a^n = a^{m+n}
			\]
		\end{block}
		
		\begin{exampleblock}{Ejemplo}
			\[
			2^3\cdot2^4 = 2^{3+4}=2^7=128
			\]
		\end{exampleblock}
		
		\begin{alertblock}{Interpretación}
			Si la base es la misma, \destacar{se suman los exponentes}.
		\end{alertblock}
	\end{frame}
	
	% ----------------------------------------------------------
	\begin{frame}{Propiedad 2: cociente de potencias de igual base}
		\begin{block}{Regla}
			\[
			\frac{a^m}{a^n}=a^{m-n}, \qquad a\neq 0
			\]
		\end{block}
		
		\begin{exampleblock}{Ejemplo}
			\[
			\frac{5^6}{5^2}=5^{6-2}=5^4=625
			\]
		\end{exampleblock}
		
		\begin{alertblock}{Interpretación}
			Si la base es la misma, \destacar{se restan los exponentes}.
		\end{alertblock}
	\end{frame}
	
	% ----------------------------------------------------------
	\begin{frame}{Propiedad 3: potencia de una potencia}
		\begin{block}{Regla}
			\[
			(a^m)^n = a^{m\cdot n}
			\]
		\end{block}
		
		\begin{exampleblock}{Ejemplo}
			\[
			(3^2)^4 = 3^{2\cdot4}=3^8=6561
			\]
		\end{exampleblock}
		
		\begin{alertblock}{Interpretación}
			En una potencia elevada a otra potencia, \destacar{se multiplican los exponentes}.
		\end{alertblock}
	\end{frame}
	
	% ----------------------------------------------------------
	\begin{frame}{Propiedad 4: potencia de un producto}
		\begin{block}{Regla}
			\[
			(ab)^n = a^n b^n
			\]
		\end{block}
		
		\begin{exampleblock}{Ejemplo}
			\[
			(2\cdot5)^3 = 2^3\cdot5^3 = 8\cdot125 = 1000
			\]
		\end{exampleblock}
		
		\begin{alertblock}{Interpretación}
			La potencia puede distribuirse a cada factor del producto.
		\end{alertblock}
	\end{frame}
	
	% ----------------------------------------------------------
	\begin{frame}{Propiedad 5: potencia de un cociente}
		\begin{block}{Regla}
			\[
			\left(\frac{a}{b}\right)^n = \frac{a^n}{b^n}, \qquad b\neq 0
			\]
		\end{block}
		
		\begin{exampleblock}{Ejemplo}
			\[
			\left(\frac{2}{3}\right)^3 = \frac{2^3}{3^3}=\frac{8}{27}
			\]
		\end{exampleblock}
		
		\begin{alertblock}{Interpretación}
			La potencia también puede distribuirse al numerador y al denominador.
		\end{alertblock}
	\end{frame}
	
	% ----------------------------------------------------------
	\begin{frame}{Propiedad 6: exponente cero}
		\begin{block}{Regla}
			\[
			a^0=1,\qquad a\neq0
			\]
		\end{block}
		
		\begin{exampleblock}{Ejemplo}
			\[
			7^0=1
			\]
		\end{exampleblock}
		
		\begin{alertblock}{Observación}
			Toda base distinta de cero elevada a cero vale 1.
		\end{alertblock}
	\end{frame}
	
	% ----------------------------------------------------------
	\begin{frame}{Propiedad 7: exponente negativo}
		\begin{block}{Regla}
			\[
			a^{-n}=\frac{1}{a^n},\qquad a\neq0
			\]
		\end{block}
		
		\begin{exampleblock}{Ejemplo}
			\[
			2^{-3}=\frac{1}{2^3}=\frac{1}{8}
			\]
		\end{exampleblock}
		
		\begin{alertblock}{Interpretación}
			Un exponente negativo indica el inverso de la potencia positiva.
		\end{alertblock}
	\end{frame}
	
	% ----------------------------------------------------------
	\begin{frame}{Propiedad 8: exponente uno y casos especiales}
		\begin{block}{Reglas}
			\[
			a^1=a
			\qquad\qquad
			1^n=1
			\qquad\qquad
			0^n=0 \text{ para } n>0
			\]
		\end{block}
		
		\begin{exampleblock}{Ejemplos}
			\[
			9^1=9
			\qquad
			1^{12}=1
			\qquad
			0^5=0
			\]
		\end{exampleblock}
	\end{frame}
	
	% ----------------------------------------------------------
	\begin{frame}{Errores frecuentes}
		\begin{alertblock}{Error frecuente 1}
			\[
			a^m+a^n \neq a^{m+n}
			\]
		\end{alertblock}
		
		\begin{exampleblock}{Ejemplo}
			\[
			2^2+2^3=4+8=12
			\]
			pero
			\[
			2^{2+3}=2^5=32
			\]
		\end{exampleblock}
		
		\pp
		
		\begin{alertblock}{Error frecuente 2}
			\[
			(a+b)^n \neq a^n+b^n
			\]
		\end{alertblock}
	\end{frame}
	
	% ----------------------------------------------------------
	\begin{frame}{Otro error frecuente}
		\begin{exampleblock}{Ejemplo comparativo}
			\[
			(2+3)^2 = 5^2 = 25
			\]
			pero
			\[
			2^2+3^2=4+9=13
			\]
		\end{exampleblock}
		
		\begin{ideaBox}
			No toda operación con potencias permite “repartir” el exponente.
		\end{ideaBox}
	\end{frame}
	
	% ----------------------------------------------------------
	\begin{frame}{Patrones numéricos}
		\begin{block}{Potencias de 2}
			\[
			2^1=2,\quad 2^2=4,\quad 2^3=8,\quad 2^4=16,\quad 2^5=32
			\]
		\end{block}
		
		\pp
		
		\begin{block}{Potencias de 10}
			\[
			10^1=10,\quad 10^2=100,\quad 10^3=1000,\quad 10^4=10000
			\]
		\end{block}
		
		\pp
		
		\begin{ideaBox}
			Cuando el exponente aumenta en 1, la potencia crece multiplicando por la base.
		\end{ideaBox}
	\end{frame}
	
	% ----------------------------------------------------------
	\begin{frame}{Práctica guiada 1}
		\begin{block}{Ejercicio}
			Simplificar:
			\[
			2^4\cdot2^3
			\]
		\end{block}
		
		\pp
		
		\begin{exampleblock}{Paso 1}
			La base es la misma, por lo tanto sumamos exponentes:
			\[
			2^4\cdot2^3 = 2^{4+3}
			\]
		\end{exampleblock}
		
		\pp
		
		\begin{exampleblock}{Paso 2}
			\[
			2^{4+3}=2^7=128
			\]
		\end{exampleblock}
	\end{frame}
	
	% ----------------------------------------------------------
	\begin{frame}{Práctica guiada 2}
		\begin{block}{Ejercicio}
			Resolver:
			\[
			(3^2)^3
			\]
		\end{block}
		
		\pp
		
		\begin{exampleblock}{Paso 1}
			Es una potencia de una potencia, así que multiplicamos exponentes:
			\[
			(3^2)^3=3^{2\cdot3}
			\]
		\end{exampleblock}
		
		\pp
		
		\begin{exampleblock}{Paso 2}
			\[
			3^{2\cdot3}=3^6=729
			\]
		\end{exampleblock}
	\end{frame}
	
	% ----------------------------------------------------------
	\begin{frame}{Práctica guiada 3}
		\begin{block}{Ejercicio contextualizado}
			En una parcela experimental se utilizan:
			\[
			2^5
			\]
			semillas por metro cuadrado.
			
			Si la superficie cultivada es de:
			\[
			2^3
			\]
			metros cuadrados, determine la cantidad total de semillas.
		\end{block}
	\end{frame}
	
	% ----------------------------------------------------------
	\begin{frame}{Resolución práctica guiada 3}
		\begin{exampleblock}{Paso 1}
			Multiplicamos las cantidades:
			\[
			2^5\cdot2^3
			\]
		\end{exampleblock}
		
		\pp
		
		\begin{exampleblock}{Paso 2}
			Como la base es la misma:
			\[
			2^5\cdot2^3=2^{5+3}=2^8
			\]
		\end{exampleblock}
		
		\pp
		
		\begin{exampleblock}{Paso 3}
			\[
			2^8=256
			\]
			Se utilizan \textbf{256 semillas}.
		\end{exampleblock}
	\end{frame}
	
	% ----------------------------------------------------------
	\begin{frame}{Aplicaciones en Agronomía}
		\begin{enumerate}
			\item \destacar{Crecimiento bacteriano:} duplicación de microorganismos en el suelo.
			\item \destacar{Semillas por hectárea:} cálculo de densidad de siembra.
			\item \destacar{Concentración de nutrientes:} uso de potencias de 10.
			\item \destacar{Escalas de medida:} cantidades muy grandes o muy pequeñas.
		\end{enumerate}
		
		\pp
		
		\begin{ideaBox}
			Las potencias permiten modelar fenómenos agrícolas de manera simple y precisa.
		\end{ideaBox}
	\end{frame}
	
	% ----------------------------------------------------------
	\begin{frame}{Aplicación contextualizada: fertilización}
		\begin{block}{Situación}
			Un fertilizante se aplica en dosis de:
			\[
			5\times10^2 \text{ gramos por hectárea}
			\]
			Si el ensayo abarca:
			\[
			10^3 \text{ hectáreas}
			\]
			¿cuántos gramos se necesitan en total?
		\end{block}
		
		\pp
		
		\begin{exampleblock}{Resolución}
			\[
			(5\times10^2)(10^3)=5\times10^5
			\]
			\[
			5\times10^5=500000
			\]
			Se requieren \textbf{500000 gramos}.
		\end{exampleblock}
	\end{frame}
	
	% ----------------------------------------------------------
	\begin{frame}{Práctica autónoma}
		Resuelva individualmente:
		\begin{enumerate}
			\item \(2^5\cdot2^2\)
			\item \(\dfrac{3^7}{3^4}\)
			\item \((4^2)^3\)
			\item \((2\cdot3)^2\)
			\item \(\left(\dfrac{3}{5}\right)^2\)
			\item \(8^0\)
			\item \(2^{-4}\)
			\item En un cultivo hay \(3^4\) plantas por sector y \(3^2\) sectores. ¿Cuántas plantas hay en total?
		\end{enumerate}
	\end{frame}
	
	% ----------------------------------------------------------
	\begin{frame}{Práctica autónoma con respuestas}
		\begin{multicols}{2}
			\begin{enumerate}
				\item \(2^5\cdot2^2=2^7=128\)
				\item \(\dfrac{3^7}{3^4}=3^3=27\)
				\item \((4^2)^3=4^6=4096\)
				\item \((2\cdot3)^2=2^2\cdot3^2=36\)
				\item \(\left(\dfrac{3}{5}\right)^2=\dfrac{9}{25}\)
				\item \(8^0=1\)
				\item \(2^{-4}=\dfrac{1}{16}\)
				\item \(3^4\cdot3^2=3^6=729\)
			\end{enumerate}
		\end{multicols}
	\end{frame}
	
	% ----------------------------------------------------------
	\begin{frame}{Cierre: síntesis}
		\begin{block}{Ideas clave de la clase}
			\begin{itemize}
				\item Una potencia representa una multiplicación repetida.
				\item La base es el número que se repite.
				\item El exponente indica cuántas veces se repite.
				\item Las propiedades permiten simplificar cálculos.
				\item Las potencias tienen aplicaciones directas en problemas de Agronomía.
			\end{itemize}
		\end{block}
	\end{frame}
	
	% ----------------------------------------------------------
	\begin{frame}{Cierre reflexivo}
		\begin{exampleblock}{Pregunta de salida}
			¿Por qué las potencias permiten representar de manera eficiente fenómenos de crecimiento en contextos agrícolas?
		\end{exampleblock}
		
		\pp
		
		\begin{ideaBox}
			\centering
			\textbf{Aprender potencias no solo ayuda a calcular: ayuda a interpretar fenómenos reales y a tomar mejores decisiones técnicas.}
		\end{ideaBox}
	\end{frame}
	
\end{document}